\documentclass[11pt]{article}
\usepackage[margin=1in]{geometry}
\usepackage{amsmath,amssymb,amsthm,enumitem,hyperref}
\hypersetup{colorlinks=true,linkcolor=blue,urlcolor=blue,citecolor=blue}
\newtheorem{theorem}{Theorem}
\newtheorem{lemma}{Lemma}
\newcommand{\MM}{\mathcal M}
\title{Branched simplicial fans satisfy the gradient--Laplacian identity\\\large MathDB 350676}
\author{Principia Math In Progress}
\date{September 22, 2026}
\begin{document}
\maketitle
\begin{abstract}
The complex with facets $012,013,014$ is not a pseudomanifold,
but its top boundary Laplacian and discrete-gradient generating
polynomial both equal $\lambda^4(\lambda+2)^2(\lambda+5)$.
This disproves the characterization in Conjecture 1 of the current
Contreras--Tawfeek preprint. We give a uniform formula for any
number of simplices joined along a common facet. The identified
journal final remains unread, so campaign eligibility is pending.
\end{abstract}

\section*{Task 1. Full counterexample to the preprint formulation}
Follow the source's conventions \cite[pp.~3,8,18]{CT}:
$B=\partial_d$ has rows indexed by the $(d-1)$-faces, columns
indexed by the $d$-faces, and $\Delta_d=BB^{\mathsf T}$.
Put $m=|K_d|$ and $n=|K_{d-1}|$.
The set $\MM_{m-i}(K)$ consists of discrete gradients with
$i$ matched pairs, all between these two dimensions; every cell
of dimension below $d-1$ is critical. Thus it is the set of
acyclic matchings of size $i$ in the top two levels of the Hasse
diagram. As in the source, $\MM_j$ is empty for $j<0$.
Conjecture 1 asserts that
\[
 \det(\Delta_d+\lambda I_n)
 =\sum_{i=0}^n |\MM_{m-i}(K)|\lambda^{n-i}                 \tag{1}
\]
holds if and only if $K$ is an orientable pseudomanifold.

\begin{theorem}
Let $d\ge2$, $r\ge3$, and let $S$ consist of $d$ vertices.
Choose distinct further vertices $u_1,\ldots,u_r$, and let
$K(d,r)$ be the simplicial complex generated by
\[
 S\cup\{u_1\},\ldots,S\cup\{u_r\}.
\]
This is a finite, connected, pure $d$-dimensional complex whose
facet-adjacency graph is connected. It is not a pseudomanifold,
and nevertheless it satisfies (1). Both sides are
\[
 \lambda^{1+r(d-1)}(\lambda+d)^{r-1}(\lambda+d+r).           \tag{2}
\]
\end{theorem}

\begin{proof}
Any two top faces meet in the common $(d-1)$-face $S$.
Each top face has $d$ other boundary faces, and each of these
belongs to that top face alone. Hence $m=r$ and $n=1+rd$.
In particular, $S$ belongs to $r\ge3$ top faces, violating
the source's requirement of exactly two. It also violates the
at-most-two condition if pseudomanifolds with boundary are
allowed. The obstruction is therefore branching, rather than
just the presence of a boundary or a disconnected complex.

Orient the top faces so that their incidence with $S$ is $+1$.
The supports of any two columns of $B$ intersect only in that
row. Every column has $d+1$ entries of absolute value $1$.
Consequently
\[
 B^{\mathsf T}B=dI_r+J_r,
\]
where $J_r$ is the all-ones matrix. Its eigenvalues are
$d+r$ once and $d$ with multiplicity $r-1$.
The nonzero eigenvalues of $BB^{\mathsf T}$ are the same,
and its nullity is $n-r$. This proves (2) for the determinant.

The bipartite incidence graph between the top two levels is a
tree: one central vertex $S$ is joined to $r$ top-face vertices,
each with $d$ additional leaves. Every matching is acyclic,
since a directed cycle would give an undirected cycle in this
tree. The remaining Hasse edges point downward and cannot
belong to a directed cycle, since no lower cell is matched.
Thus every such matching is a discrete gradient of the
required type. For completeness, this last assertion can be
seen without an existence theorem: assign strictly decreasing
positive values along the arrows in an acyclic orientation of
these two levels, and assign smaller, dimension-increasing
values to the lower levels. The resulting function has exactly
the matched pairs as its reversed inequalities and satisfies
the defining discrete Morse inequalities.

Choose the $i$ top faces that will be matched. If the central
face $S$ is unused, there are $d^i$ choices. If it is used,
there are $i$ choices for its partner and $d^{i-1}$ choices
for the others. Therefore, for $1\le i\le r$,
\[
 |\MM_{r-i}(K(d,r))|=\binom ri(d^i+i d^{i-1}),             \tag{3}
\]
and the empty matching gives $|\MM_r|=1$. The binomial theorem
and its derivative now give
\begin{align*}
 \sum_{i=0}^n |\MM_{r-i}|\lambda^{n-i}
 &=\lambda^{n-r}\bigl((\lambda+d)^r
                   +r(\lambda+d)^{r-1}\bigr)\\
 &=\lambda^{n-r}(\lambda+d)^{r-1}(\lambda+d+r).
\end{align*}
This is (2), completing the proof.
\end{proof}

\subsection*{The explicit five-vertex example}
Take $S=\{0,1\}$ and $u_1=2,u_2=3,u_3=4$.
For rows $01,02,03,04,12,13,14$ and columns $012,013,014$,
the usual increasing-vertex orientations give
\[
 B=\begin{pmatrix}
 1&1&1\\ -1&0&0\\0&-1&0\\0&0&-1\\
 1&0&0\\0&1&0\\0&0&1
 \end{pmatrix},\qquad
 (|\MM_3|,|\MM_2|,|\MM_1|,|\MM_0|)=(1,9,24,20).
\]
Both polynomials are
\[
 \lambda^7+9\lambda^6+24\lambda^5+20\lambda^4.
\]
The equivalence is therefore false. This counterexample does
not refute its separate forward implication for orientable
pseudomanifolds.

\subsection*{Why the identity also holds for incidence forests}
More generally, (1) holds whenever the incidence graph of the
top two levels is a forest. Here is a direct matrix explanation.
Expansion in principal minors followed by Cauchy--Binet gives
\[
 \det(\lambda I_n+BB^{\mathsf T})
 =\sum_{i=0}^{\min(m,n)}\lambda^{n-i}
   \sum_{\substack{U\subseteq[n],\ V\subseteq[m]\\|U|=|V|=i}}
       \det(B_{U,V})^2.                                    \tag{4}
\]
Two different perfect matchings on the same selected vertices
would have an alternating cycle in their union. Thus an
incidence forest has at most one such matching. Each nonzero
determinant in (4) is consequently $1$ or $-1$, and its square
counts that unique matching. All matchings are acyclic, so (4)
is exactly (1). This is an elementary application of
Cauchy--Binet and matching enumeration, not a claim that those
classical tools are new.

\section*{Tasks 2--5. Proposed follow-up problems}
\begin{enumerate}[start=2,leftmargin=*]
\item Classify the complexes satisfying (1) when their top
incidence graph is connected and has exactly one cycle,
in terms of its signed incidences.
\item Determine whether (1) admits a pure, strongly connected
two-dimensional counterexample to the characterization in which
every edge belongs to at least two triangles.
\item Determine which stellar subdivisions preserve (1),
starting with subdivisions of the common face in $K(d,r)$.
\item For fixed $d$ and $r$, determine the largest number of
top-dimensional discrete gradients among pure connected
$d$-complexes with $r$ facets and a forest as their top
incidence graph.
\end{enumerate}

\paragraph{Verification and review.}
The accompanying exact-arithmetic script constructs boundary
matrices from simplex lists, enumerates matchings and checks
their directed acyclicity, and compares their counts with
the characteristic polynomial. The proof above establishes
the whole family. All review was by the authoring agent;
no independent peer review or proof-assistant certification
is claimed.

\begin{thebibliography}{9}
\bibitem{CT} I.~Contreras and A.~R.~Tawfeek,
\emph{On discrete gradient vector fields and Laplacians of
simplicial complexes},
\href{https://arxiv.org/abs/2105.05388v2}{arXiv:2105.05388v2},
2 July 2022, 20 pages; Conjecture 1, p.~18.
An identified journal version appeared in \emph{Annals of
Combinatorics} \textbf{28} (2024), 67--91,
\href{https://doi.org/10.1007/s00026-023-00655-1}{doi:10.1007/s00026-023-00655-1}.
The journal version has not been read in this audit.
\end{thebibliography}
\end{document}
